\section{Tempo e spazio}

\subsection{Complessitá di tempo}
$M(x)$ richiede tempo $t$ sse:
\[
\configurazionestart{x} \xrightarrow{M^t} \configurazione{q_H}{w}{u}
\]
Dove $q_H \in \{h, "yes", "no"\}. $
\begin{itemize}
	\item {Una MdT deterministica $M$ {opera} in tempo $f(n)$ sse $\forall x \in \Sigma^*$ $M(x)$ richiede tempo $t \leq f(n)$} 
	\item {Una MdT non deterministica $N$ {opera} in tempo $f(n)$ sse $\forall x \in \Sigma^*$ tutti i run della computazione di $N(x)$ terminano in tempo $t \leq f(n)$} 
\end{itemize}


\subsection{Complessitá di spazio}
Si potrebbe considerare la somma della lunghezza dei $k$ nastri: $\sum_{i=1}^{k} |w_i u_i|$ \\
Ma possiamo scegliere anche la lunghezza del nastro piú lungo: $\max_{i \in \{1, \dots, k\}} |w_i u_i|$. \\
Queste due quantitá sono asintoticamente equivalenti:\\
$\max_{i \in \{1, \dots, k\}} |w_i u_i| \leq  \sum_{i=1}^{k} |w_i u_i| \leq k \cdot \max_{i \in \{1, \dots, k\}} |w_i u_i|$  \\

\textbf{Attenzione:} con queste definizioni naive, lo spazio occupato dall'input verrá sempre contato.
Questo vuol dire che lo spazio é sempre limitato inferiormente da $n$, anche quando non usiamo variabili ausiliarie. \\
A noi interessa considerare lo \textit{spazio ausiliario} utilizzato. \\
Per farlo, occorre considerare le macchine di Turing con input e output e contare solo i nastri interni.

\subsubsection{Complessitá di spazio ausiliario nelle MdT con input e output}
$M(x)$ richiede spazio $z$ sse: \\
Se:
\[
\configurazionestart{x} \xrightarrow{M^*} \configurazione{q_H}{w}{u}
\]
Dove $q_H \in \{h, "yes", "no"\} \quad$ e {$\quad z = \sum_{i=2}^{k-1} |w_i u_i|$}

\begin{itemize}
\item {Una MdT deterministica $M$ {opera} in spazio $g(n)$ sse $\forall x \in \Sigma^*$ $M(x)$ richiede spazio $z \leq g(n)$}
\item {Una MdT non deterministica $N$ {opera} in spazio $g(n)$ sse $\forall x \in \Sigma^*$ tutti i run della computazione di $N(x)$ terminano in spazio $z \leq g(n)$}
\end{itemize}


\subsection{Classi di complessitá}
\subsubsection{TIME}
Un linguaggio $L$ appartiene alla classe $\TIME(f(n))$ se esiste una macchina di Turing $M$ che decide $L$ e opera in tempo $f(n)$.
\subsubsection{NTIME}
Un linguaggio $L$ appartiene alla classe $\NTIME(f(n))$ se esiste una macchina di Turing non deterministica $N$ che decide $L$ e opera in tempo $f(n)$.
\subsubsection{SPACE}
Un linguaggio $L$ appartiene alla classe $\SPACE(f(n))$ se esiste una macchina di Turing $M$ che decide $L$ e opera in spazio $f(n)$.
\subsubsection{NSPACE}
Un linguaggio $L$ appartiene alla classe $\NSPACE(f(n))$ se esiste una macchina di Turing non deterministica $N$ che decide $L$ e opera in spazio $f(n)$.
\subsubsection{P}
$\P = \bigcup_{j>0} \TIME(n^j)$
\subsubsection{NP}
$\NP = \bigcup_{j>0} \NTIME(n^j)$
\subsubsection{EXP}
$\EXP = \bigcup_{j>0} \TIME(2^{n^j})$
\subsubsection{NEXP}
$\NEXP = \bigcup_{j>0} \NTIME(2^{n^j})$
\subsubsection{L}
$\L = \bigcup_{j>0} \SPACE(\log(n)^j)$
\subsubsection{PSPACE}
$\PSPACE = \bigcup_{j>0} \SPACE(n^j)$

\subsection{Classi di complessitá complementari}
Per ogni classe di complessitá $C$ definiamo la classe complementare $\co C$ come l'insieme dei linguaggi complementari a quelli di $C$. \\
$\co C = \{ \overline{L} : L \in C \}$


\subsection{Teorema: il tempo é limitato inferiormente dall'input}
{$f(n) = \Omega(n)$} \\
O, equivalentemente: \\
{$n = O(f(n))$}
\begin{proof}
Se la codifica dell'input é \textit{ragionevole}, allora la macchina di Turing deve leggere tutto l'input prima di poter decidere se accettarlo o rifiutarlo.
\end{proof}


\subsection{Teorema: lo spazio é limitato superiormente dal tempo}
{$g(n) = O(f(n))$}
\begin{proof}
Supponiamo che una macchina $M$ a 1 nastro operi in tempo $f(n)$ e che a ogni passo scriva sul nastro e sposti la testina a destra. \\
Allora, dopo $f(n)$ passi, la testina avrá scritto esattamente $g(n) = f(n)$ simboli sul nastro. \\
Qualsiasi altra macchina che operi in tempo $f(n)$ avrá eseguito dei passi sovrascrivendo celle giá utilizzate. In questo caso l'utilizzo di spazio é inferiore a quello di $M$.
\end{proof}
\subsubsection{Corollario: $\TIME(f(n)) \subseteq \SPACE(f(n))$}
\subsubsection{Corollario: $\P \subseteq \PSPACE$}
Qualunque linguaggio che puó essere deciso in tempo $f(n)$ puó essere deciso anche in spazio $f(n)$.


\subsection{Teorema: Terminazione delle MdT space bounded}
Se una MdT $M$ opera in spazio $g(n)$, allora esiste una MdT $M'$ che opera in spazio $O(g(n))$ e termina per ogni input entro $O(c^{g(n)})$ passi.
\footnote{Questo teorema presuppone il teorema dello speedup multinastro (\ref{teorema:speedup_multinastro})}
\begin{proof}
$M' = (M_1, M_2)$ é la concatenazione di due macchine, e utilizza un nastro aggiuntivo. \\
$M_1$ scrive sul nastro aggiuntivo una stringa che codifica la configurazione di $M$, fatta in questo modo: \\
\[[\triangleright, q, \sqcap, \dots, \sqcap, 1, \dots, 1]\]
\begin{itemize}
	\item Il primo simbolo codifica lo stato corrente della macchina $M$ utilizzando $|Q|$ nuovi simboli.
	\item I successivi $g(n)$ simboli $\sqcap$ codificano il contenuto del nastro di $M$, utilizzando i $|\Sigma|$ simboli dell'alfabeto di $M$.
	\item Gli ultimi $g(n)$ simboli codificano la posizione della testina di $M$, rappresentata in unario tra le sue $g(n)$ posizioni possibili.
\end{itemize}
La macchina $M_1$ é simile a $M_g$, perché rappresenta due volte $g(n)$ in unario: una prima volta con $\sqcap^{g(n)}$ e una seconda con $1^{g(n)}$. \\
Lo spazio utilizzato é $O(g(n))$. \\

$M_2$ legge la configurazione di $M$ e utilizza gli altri nastri per simularla. \\
Il nastro con la configurazione viene trattato come un contatore, che viene incrementato dopo ogni simulazione. \\
Occorre solo fare attenzione al fatto che la base non é sempre la stessa:
\begin{itemize}
	\item La prima cella del contatore puó assumere $|Q|$ valori.
	\item Le successive $g(n)$ celle del contatore possono assumere $|\Sigma|$ valori.
	\item Le ultime $g(n)$ celle del contatore sono in unario.
\end{itemize}
Se la simulazione di $M$ termina, allora anche $M_2$ termina. \\
Se invece la simulazione non termina, il contatore prima o poi esaurirá le configurazioni possibili: a quel punto, faremo terminare $M_2$ anche se $M$ non termina. \\
Il massimo numero di configurazioni possibili é $|Q| \cdot |\Sigma|^{g(n)} \cdot {g(n)} = O(c^{g(n)})$.
\end{proof}
Questo teorema ci mostra come forzare la terminazione di una MdT che opera in spazio limitato.
\subsubsection{Corollario: $\SPACE(g(n)) \subseteq \TIME(c^{g(n)})$}
\subsubsection{Corollario: $\PSPACE \subseteq \EXP$}

\subsection{Teorema: $\TIME(f(n)) \subseteq \NTIME(f(n))$}
\begin{proof}
Le macchine di Turing deterministiche sono un caso particolare di quelle non deterministiche, in cui la relazione di transizione $\Delta$ associa a ogni coppia in $Q \times \Sigma$ una sola tripla di $Q \times \Sigma \times \{\leftarrow, \rightarrow, -\}$.
\end{proof}
\subsubsection{Corollario: $\P \subseteq \NP$}
\subsubsection{Corollario: $\EXP \subseteq \NEXP$}


\subsection{Teorema: $\NTIME(f(n)) \subseteq \bigcup_{c>1} \TIME(c^{f(n)})$}
\label{theorem:non_deterministic_to_deterministic}
Data una MdT non deterministica $N$ che decide $L$ e opera in tempo $f(n)$. \\
Allora esiste una MdT deterministica $M$ che decide $L$ e opera in tempo $O(c^{f(n)})$.
\begin{figure}[h!]
	\caption{\small Albero di computazione di una MdT non deterministica.}
	\centering
	\label{fig:albero_computazione}
	\includegraphics[width=0.8\textwidth]{albero_computazione.png}
\end{figure}
\begin{proof}
	Simuliamo $N$ con una MdT deterministica $M$. \\
	Possiamo vedere la computazione di $N$ come un albero di altezza $f(n)$ (figura \ref{fig:albero_computazione}). \\
	Ogni passo corrisponde a una \textit{scelta} tra al piú $d$ possibilitá. \\
	Ogni sequenza di $t$ passi corrisponde a una sequenza di $t$ scelte. \\
	Codifichiamo quindi un run con una stringa di $t$ numeri da $1$ a $d$. \\
	$M$ tiene su un nastro la stringa che codifica una sequenza di $t$ scelte. \\
	Per una sequenza di $t$ passi, $M$ simula le transizioni che avrebbe fatto $N$. \\
	Poi incrementa la stringa e ripete la simulazione. \\
	In questo modo esplora tutti i possibili run di $N$. \\
	$M$ ha bisogno di un altro nastro per conservare sempre l'input $x$ e altri $k$ nastri per simulare quelli di $N$. \\
	Se la simulazione di un run termina in $"yes"$, $M$ accetta. \\
	Se tutte le simulazioni terminano in $"no"$, $M$ rifiuta. \\
	Il numero di possibili run é $\sum_{i=1}^{f(n)} d^i $. \\	
	Il tempo di una simulazione é $f(n)$, cioé quello che impiega $N$. \\
	$f(n) \cdot \sum_{i=1}^{f(n)} d^i = O(d^{f(n)})$.
\end{proof}
Vedremo che invece il guadagno in spazio di una macchina non deterministica rispetto a una deterministica é limitato.
\subsubsection{Corollario: $\NP \subseteq \EXP$}


\subsection{Proposizione: chiusura delle classi di complessitá deterministiche rispetto al complemento}
Se $C$ é una classe di complessitá deterministica, allora $\co C = C$
\begin{proof}
Sia $L \in C$ e sia $M$ una macchina di Turing che decide $L$ e opera in tempo $f(n)$ e spazio $g(n)$. \\
Allora possiamo costruire una macchina $M'$ che decide $\overline{L}$ e opera in tempo $f(n)$ e spazio $g(n)$, semplicemente invertendo gli stati $"yes"$ e $"no"$  di $M$. \\
\end{proof}

\subsubsection*{Perché non é possibile fare lo stesso ragionamento per le classi non deterministiche?}
Con una macchina non deterministica $N$ non basta invertire gli stati terminali. \\
Quello che vorremmo é una macchina $N'$ che accetta $\overline{L}$ perché:
\begin{itemize}
\item $x \in \overline{L}$ sse $\exists$ un run di $N'$ che accetta $x$.
\item $x \notin \overline{L}$ sse $\forall$ run di $N'$, il run non accetta $x$.
\end{itemize}
Invertendo gli stati $"yes$ e $"no"$ di $N$ non si otterrebbe una macchina che accetta il complemento del linguaggio. \\ 
Invertire $"yes"$ e $"no"$ porta una macchina che accetta $x$ se un solo run di $N$ lo rifiutava. 
\textbf{Un input accettato da una macchina non deterministica puó essere rifiutato da un run}, e quindi accettato anche dalla nuova macchina.

\subsubsection*{}
Vedremo che é possibile dimostrare la chiusura delle classi di complessitá di spazio non deterministiche rispetto al complemento. \\
Per quanto riguarda le classi di complessitá di tempo non deterministiche, invece, il problema è aperto. \\

\subsection{Domande aperte}
\begin{itemize}
    \item $\P = \NP$?
    \item $\NP = \co\NP$?
    \item $\NP = \PSPACE$?
\end{itemize}
